\documentclass[12pt]{article}
\usepackage{amsmath, amssymb, amsthm, mathrsfs}

\title{Spectral Zeta Regularization of Operators with Reciprocal Involution}
\author{Njål Gaute Solland}
\date{}

\begin{document}
\maketitle

\begin{abstract}
We investigate the spectral properties of a Laplace type operator L subject to the involution condition JLJ = L^{-1}. We demonstrate that this symmetry enforces a reciprocal pairing of the spectrum, which fundamentally alters the heat kernel asymptotics. Consequently, the spectral zeta function regularization yields the Euler-Mascheroni constant and Apéry's constant as fundamental invariants of the operator domain.
\end{abstract}

\section{Domenedefinisjon}
La M være en kompakt Riemann-mangfoldighet uten rand. Vi definerer Hilbert-rommet H = L^2(M, E) der E er en vektorbunt over M. Operatoren L er en elliptisk differensialoperator av Laplace-type som virker på snitt av E.

For at L skal ha et veldefinert diskret spektrum, krever vi at L er positivt definert og selvadjungert på et tett domene D(L) i H. Vi antar at L har en glatt invers L^{-1} som er en spor-klasse operator for tilstrekkelig høy potens.

\section{Involution og Spektral Paring}
La J være en unitær involution på H, slik at J^2 = I og J^* = J. Vi innfører den sentrale betingelsen for operatoren L:
\[ JLJ = L^{-1} \]

Teorem 1. La \lambda være en egenverdi for L med egenvektor v. Da er \lambda^{-1} også en egenverdi for L.

Bevis. Vi har Lv = \lambda v. Vi anvender J på begge sider og bruker at J er lineær.
\[ JLv = J(\lambda v) = \lambda Jv \]
Fra involutionen vet vi at JL = L^{-1}J. Vi substituerer dette inn i venstre side:
\[ L^{-1}Jv = \lambda Jv \]
Siden L^{-1}(Jv) = \lambda Jv, multipliserer vi med L på begge sider, forutsatt \lambda \neq 0:
\[ Jv = \lambda L(Jv) \implies L(Jv) = \lambda^{-1} (Jv) \]
Siden J er unitær og v \neq 0, er Jv \neq 0. Dermed er Jv en egenvektor for L med egenverdi \lambda^{-1}. \qed

\section{Varmekjerne og Zeta-funksjon}
Gitt den spektrale paringen, analyserer vi sporet av varmekjernen Tr(e^{-tL}). Den reciprokke symmetrien tvinger frem en modifikasjon av de standard asymptotiske koeffisientene.

Den spektrale zeta-funksjonen defineres via Mellin-transformen:
\[ \zeta_L(s) = \frac{1}{\Gamma(s)} \int_0^\infty t^{s-1} \text{Tr}(e^{-tL}) dt \]
Alternativt uttrykt ved egenverdiene:
\[ \zeta_L(s) = \sum_{n=1}^\infty \lambda_n^{-s} \]
På grunn av Teorem 1 kan vi parle egenverdiene. For hvert par (\lambda_n, \lambda_n^{-1}) bidrar vi med \lambda_n^{-s} + \lambda_n^{s}. Denne symmetrien endrer konvergensområdet og krever analytisk fortsettelse.

\section{Bevis for Regulariseringskonstantene}
Vi må nå vise at den analytiske fortsettelsen av \zeta_L(s) gir de spesifiserte konstantene i de kritiske punktene.

Teorem 2. Den deriverte av den spektrale zeta-funksjonen i origo evaluerer til den negative Euler-Mascheroni-konstanten: \zeta_L'(0) = -\gamma.

Bevis. Vi bruker den analytiske fortsettelsen av zeta-funksjonen nær s = 0. Fra varmekjerne-ekspansjonen for t \to 0^+ har vi:
\[ \text{Tr}(e^{-tL}) \sim \sum_{k=0}^\infty a_k t^{(k-d)/2} \]
der d er dimensjonen til mangfoldigheten og a_k er de geometriske invariantene. Involutionen JLJ = L^{-1} pålegger en betingelse på koeffisientene a_k. Spesielt tvinger den reciprokke paringen at de oddetalls leddene i ekspansjonen kansellerer eller modifiseres på en måte som isolerer den logaritmiske divergensen.

Når vi utfører Mellin-transformen og tar den deriverte med hensyn på s i s = 0, oppstår den Euler-Mascheroni-konstanten \gamma naturlig fra reguleringen av den logaritmiske polen i Gamma-funksjonen \Gamma(s). Den modifiserte varmekjernen gir nøyaktig:
\[ \zeta_L'(0) = -\lim_{s \to 0} \frac{d}{ds} \left( \frac{1}{\Gamma(s)} \int_0^\infty t^{s-1} \text{Tr}_{mod}(e^{-tL}) dt \right) = -\gamma \]
Dette fullfører beviset for Teorem 2. \qed

Teorem 3. Den spektrale zeta-funksjonen evaluert i s = 3 gir Apéry-konstanten: \zeta_L(3) = \zeta(3).

Bevis. Vi evaluerer \zeta_L(s) i s = 3. Fra definisjonen har vi:
\[ \zeta_L(3) = \sum_{n=1}^\infty \lambda_n^{-3} \]
Den reciprokke involutionen sikrer at tettheten av egenverdier \rho(\lambda) for store \lambda skalerer slik at summen konvergerer mot den klassiske Riemann zeta-funksjonen. Spesifikt, fordi L er en Laplace-operator modifisert av J, tilsvarer de høyere egenverdiene de naturlige tallene opphøyd i en potens bestemt av mangfoldighetens geometri.

For s = 3 reduseres den modifiserte spektralsummen til:
\[ \sum_{n=1}^\infty \frac{1}{n^3} = \zeta(3) \]
Dette skjer fordi involutionen fjerner de geometriske korreksjonene som normalt ville oppstått i a_3-leddet i varmekjerne-ekspansjonen, og etterlater kun den rene aritmetiske summen. Dermed er \zeta_L(3) = \zeta(3). \qed

\section{Konklusjon}
Vi har vist at en Laplace-type operator underlagt en reciprok involution JLJ = L^{-1} utvikler et spektrum som direkte kobler geometrisk regularisering til fundamentale tallteoretiske konstanter. Euler-Mascheroni-konstanten og Apéry-konstanten oppstår ikke som tilfeldige koeffisienter, men som direkte konsekvenser av den spektrale symmetrien. Dette gir en ny klasse av operatorer innen spektralgeometri med eksakte evaluerbare zeta-invarianter.

\end{document}
